% !TEX program = xelatex
\documentclass[12pt,a4paper]{ctexart}

\usepackage{fontspec}
\setmainfont{Times New Roman}
\setCJKmainfont{SimSun}[BoldFont=SimHei,ItalicFont=KaiTi]
\setCJKsansfont{SimHei}
\setCJKmonofont{FangSong}

\usepackage[top=28mm,bottom=28mm,left=28mm,right=28mm]{geometry}
\usepackage{fancyhdr}
\usepackage{perpage}
\MakePerPage{footnote}

\setlength{\parindent}{2em}
\setlength{\parskip}{0pt}

\newcommand{\circledfootnotenumber}{%
  {\fontspec{SimSun}\ifcase\value{footnote}\or ①\or ②\or ③\or ④\or ⑤\or ⑥\or ⑦\or ⑧\or ⑨\else\arabic{footnote}\fi}}
\renewcommand{\thefootnote}{\circledfootnotenumber}

\pagestyle{fancy}
\fancyhf{}
\fancyfoot[C]{\thepage}
\renewcommand{\headrulewidth}{0pt}

\begin{document}

\begin{center}
    {\zihao{-3}\heiti 半个世纪后，我想对您说：我不是一下子明白的\par}
    \vspace{0.8em}
    {\zihao{-4}\kaishu 姓名：\underline{\hspace{3.5cm}}\qquad 学号：\underline{\hspace{4cm}}\par}
\end{center}

\vspace{1em}

{\zihao{-4}\songti\setlength{\baselineskip}{20pt}\selectfont

毛泽东同志：

时光荏苒，您走了已经半个世纪了。作为一个后辈，我借着这次课程总结的机会写下这些话。它承载着我想对您说的一些话，这每一句话，都是我在这一学期的感悟。
我当然不是突然明白毛泽东思想的，而是在一个循序渐进的过程中逐渐明白的。

之前的马原思政课更像价值观的入门；而到了这门课，我以为会离我们的生活很远，其实后来发现并不然。学期之初，确实有些把毛概课应付过去的心态，因为可能之前用第一志愿抢到
的这门课，考虑到的有些因素就是给分好、没有大作业，想的是更多的心思放到专业课上去。听了几节课之后，感觉老师讲课富有饱满的热情，努力把历史讲得
活灵活现，而不是照着PPT平铺直叙地念。这让我慢慢愿意进入课程本身。

从小学习过近现代史的我们，对于很多概念都耳熟能详，比如走群众路线、要独立自主。读《新民主主义论》\footnote{毛泽东：《新民主主义论》，《毛泽东选集》第2卷，北京：人民出版
社，1991年，第662--706页。}，我才意识到如果把它们放回当时的历史处境中，其实想做到的话并不轻松。文章一开头问“中国向何处去”，初次见时，觉得它有些意
味深长。再看到时，才发自内心地觉得它沉重。那时候的中华民族，没有在几条现成道路之间选择的空间，而是在外部压迫不断加深、社会结构被撕裂的情况下，必须迫
切找出一条走得通、走得好的路。

以前的我，把道路理解得太直观了。仿佛天然就是方法和途径的意象。这也就容易产生一种错觉，好像后来被证明是对的道路，它的产生和发展就自然而然的。
这个学期听课以后，我才发觉不是这样。鲁迅说过，世上本没有路，走的人多了，也便成了路。中国革命的道路也不可能是先铺好再走，而是在失败中判断，在
碰壁中探索。主席，您提出的新民主主义，不是因为前面加了一个新字而显得新颖，而是因为当时中国的社会性质决定了革命道路绝对不能照搬别人。
在新民主主义论中，我印象最深的还是从实际出发。再到后来1978年关于真理标准问题的大讨论又强调了这一点。

课堂上，老师讲到您的历史功绩，也讲到晚年探索的曲折。过去我更习惯听到一些确定的评价，像是您作为伟大领袖，带领中国人民走向革命胜利，建立了新中国。
果只停在这里，我们好像离真正的历史还是很远。后来听到那些曲折和代价，我反而更能感觉到，历史人物不是一张扁平的图片，而是饱满的、立体的。新中国从一穷
二白、百废待兴中起步，能够逐步建立比较完整的工业体系和国防体系，这不是几句轻飘飘的话就能一笔带过的成就。可是我们探索中的曲折也提醒我们，理想如果脱离
实际，或者缺少及时纠正，代价会很大。我们不能报喜不报忧，不能只看光明，而故意不看它背后的阴影。

课程之中，有一个词被反复提到，那就是“人民”。它在我国的政治理论中出现得太多，多到有时候会让我产生一种错觉，忘记了它本来是非常具体的存在。读《湖南农民运动考察报告》时
，我才有身临其境的感觉。您不是高高在上，把人民看成远处模糊的一片，而是真正深入到田间地头，去观察他们为什么而愤怒，为什么而行动，又如何成为革命的力量。
\footnote{毛泽东：《湖南农民运动考察报告》，《毛泽东选集》第1卷，北京：人民出版社，1991年，第12--44页。}所谓群众路线，如果只是剩下“从群众中来
，到群众中去”这样一句话，那么它的价值就被磨平了。要走群众路线，首先要能看见普通人的生活，承认人民不被动等待安排，而是能够主动改变历史。

现在互联网舆论中，我们经常能看到一些情绪化的判断。大众的选择有时会先受到片面信息的影响，再被即时情绪所左右，最后又和局部利益纠缠在一起。这种现象确实存在。
但如果因为这一点就把人民看低，甚至觉得普通人只能被安排，那更危险。也许真正困难的地方就在这里，既要承认人民中有真实的力量，也要承认这种力量需要被组织、
被引导，被放到更长远的整体利益中去理解。这样再回头看人民二字，它就不可能只是一句漂亮话，而是一个真正需要党和国家来面对的现实问题。

我自己当然也是人民群众的一部分，只不过作为学生，生活还没有完全融入社会。我的生活很容易被缩小成一些指标，作业、考试、绩点、证书、升学或就业去向。这些都是
现实的问题，对我也都很重要。但如果只盯着这些指标看，人是不是也会变得扁平化，路也走窄了？甚至变成一种单向度的人？这个道路，也刚好和前面讨论过的
道路问题联系起来。读《青年运动的方向》时，我注意到文章里谈到青年和人民结合的问题。\footnote{毛泽东：《青年运动的方向》，《毛泽东选集》第2卷，北京：人
民出版社，1991年，第559--565页。}今天的条件显然完全不同了，我们不可能简单套用当年的方法，但这个问题仍然提醒着我，青年不能只在自己的象牙塔里理解世界。
一个人学什么、以后做什么，都会和社会发生关系，只是这种关系没有那么直接和明显。

以上也是我对青年责任的一点理解。以前看到“青年责任”这个词，会觉得很大，离我很远。现在我觉得，我们不一定要从特别宏大的叙事中寻找它，也可以从具体
生活的方方面面去体会。面对一些网络舆论事件时，我们也不要急于站队。很多事一旦被片面化，就很容易让我们陷入虚假两难的窘境。更重要的是先尽可能多地了解事实真相，形成自己的判断。
这其实也回到了我们之前讨论的从实际出发。看到房地产市场的起伏，或者一个专业从冷门到热门的变化，我们也不应该只停留在情绪之中，而应该去更多地了解相关背景，去思考
相关发展是在什么背景、什么条件下作出的、如何发展至今。这样说可能还是有些老套，但它总比空泛地说我要贡献力量，更接近作为学生，我们真实的状态。具体到我的学习生活中，学习专业知识的时候，也不应该只顾及考试与做题而不注重背后的原理。有些课能直接训练我们未来解决问题的能力，有些课则提高了我们的综合素质。

有时候我也会觉得，自己写这些话好像在讲官话套话，但这也许是这门课的教学形式所致。我们必须从课本里的术语和概念切入，再一点点把它们和自己的主观世界、过往经历衔接起来。
接不上时，我们可能会觉得有些空；但能勉强接上一点点时，才会发现这些上个世纪形成的思想，确实能和眼前的现实产生联系。当然，作为课程，它本身也有一定的局限。它没有让我们
把所有问题都学会想明白，有些内容我仍然觉得有些虚无缥缈，有些抽象，很多历史事件也不是在两节课内，在读一两篇文章之后就能完全了解掌握的。但它至少给了我看问题的角度。
以前我很容易先找结论再套用，现在会提醒自己先看条件，为什么当时是这样的情况，为什么不能那样，如果换到今天，这个方法还适用吗？这些问题绕来绕去，看起来有些麻烦。
但也正是这种麻烦，让我觉得思想政治理论课不是停留在简单的背诵。

半个世纪后，我想对您说，我对您的理解，对中国近现代史的理解，都谈不上深，但总比半年前具体了一些。您留给后人的不仅仅是革命胜利的历史记忆，还有一种回到现实中寻找道路的感觉。它提醒人不要把理论说得太过美好，也不要把现实看得太过简单。我们当然要有理想，但也要承认，理想需要一步步在实际中印证践行。于我而言，这就是我在这门课中的收获。

我现在还不能回答很多大的问题，也不能很清楚地看清问题的根本在哪里、制约条件有哪些。我当然不敢说自己已经完全理解中国道路是什么，但我至少
开始意识到，理解中国本身就是一个长期学习的过程。之后无论做什么，我也应该时常提醒自己，不要只把自己关在个人得失的囹圄之中。能够多理解一些历史，多看见一点
现实，把自己的选择和国家、社会联系起来，也许这就是我作为一个普通青年可以先做的事情。

\vspace{1em}
\noindent{\heiti 参考文献}

\noindent [1] 毛泽东：《新民主主义论》，《毛泽东选集》第2卷，北京：人民出版社，1991年。

\noindent [2] 毛泽东：《湖南农民运动考察报告》，《毛泽东选集》第1卷，北京：人民出版社，1991年。

\noindent [3] 毛泽东：《青年运动的方向》，《毛泽东选集》第2卷，北京：人民出版社，1991年。

}

\end{document}
